% =============================================================================
% Chapter 8 --- Limitations and Future Work
% =============================================================================

\chapter{Limitations and Future Work}
\label{ch:limitations}

This chapter catalogs the limitations of the present work and proposes
a concrete engineering roadmap for follow-on effort, grouped by their
architectural source: the boundary between the digital controller and
the analog core, the structural residue of the FSM design, and the
verification scope.

\section{Limitations and Technical Debt}
\label{sec:lim-limitations}

\paragraph{Behavioral mock vs.\ analog reality.}
The behavioral mock of \Cref{ch:verification} captures the digital
contract at the boundary --- shadow weight matrix, readback line,
\emph{op\_done} handshake --- but discards every analog observable. The
scope of what is validated is therefore explicitly \emph{control
sequencing} and \emph{digital packet semantics} against a first-order
behavioral abstraction, not full accelerator
accuracy~\cite{Ielmini2018RRAM, Sebastian2020IMC}. Three specific
caveats follow. The \emph{physical timing of} \emph{op\_done} is
testbench-generated and does not yet have a silicon-characterized
tolerance band against charge-pump settling and in-cell transients. The
\emph{end-to-end MVM numerical accuracy} cannot be claimed because the
environment has no golden MAC reference: the shadow matrix captures the
per-cell quantized weight, not the analog MVM result. The
\emph{analog non-idealities} of the array --- inter-cell cross-talk, IR
drop along long lines, retention drift between programming and verify,
and device-to-device variability beyond the discrete four-bit readback
injection patterns --- are not modeled.

\paragraph{Reset state unreachable from idle.}
The main FSM defines a reset state that asserts the core-reset output
for one cycle, but no transition from idle currently targets it. The
state and its associated reset-pulse logic are structurally present and
would behave correctly if entered. The debt is to either remove the
dead branch after an architectural audit or add a host-visible
global-reset command that drives idle into reset; the current code
absorbs either decision without re-architecting.

\paragraph{Latent signals and module parameters.}
A small number of internal signals --- a verify-initiated erase flag, an
erase-retry-exhausted flag, a weight-read enable line --- are driven by
the combinational FSM but not consumed elsewhere. They exist as
documentation hooks for a future integration (e.g.\ lifting
erase-retry-exhausted to a top-level status pin). The controller also
exposes an address-width and a weight-width parameter at the module
boundary whose values are not referenced inside the body; datapath
widths are driven by package-level constants instead. In each case the
debt is the same: either wire the hook through and add test coverage,
or remove it from the public API to avoid misleading integrators.

\paragraph{Dual-flow programming residue.}
The RTL contains two latent programming flows: the direct-address flow
used by the present testbenches, in which the captured address register
drives every per-cycle buffer access, and a legacy buffer-index-sweep
flow in which an internal counter walks a range of cells. The
direct-address flow is the path the regression exercises; the sweep
flow is partially present but is not the path the recovery logic uses.
A maintenance pass would either unify the two flows behind a single
addressing mechanism or formally deprecate the sweep.

\paragraph{Verification scope.}
The verification environment is directed. There is no UVM
infrastructure, no constrained-random stimulus engine, no functional
coverage model, and no formal property set. The absence of these
techniques is a scope decision, not a feasibility claim, and their
inclusion is the natural next step (\Cref{sec:future-sva,sec:future-cr}).

\section{Future Work --- SystemVerilog Assertions}
\label{sec:future-sva}

SystemVerilog assertions (SVA) bound to the controller as a verification
overlay would directly mitigate several limitations above. In every
case the assertion lives in a binding module that observes the relevant
internal state without modifying the synthesizable RTL. The natural
targets are the protocol boundaries that the current testbenches verify
by directed expectation, which an assertion would extend to a continuous
guard. \Cref{tab:future-sva} lists the architectural concerns and the
RTL section that grounds each one.

\begin{table}[h]
    \centering
    \caption{Architectural concerns and proposed SystemVerilog
        assertions.}
    \label{tab:future-sva}
    \begin{tabularx}{\linewidth}{X l}
        \toprule
        \textbf{Concern and proposed assertion} & \textbf{Grounding} \\
        \midrule
        \textbf{Handshake liveness.} If \emph{valid} rises while the FSM
            is in idle with a legal command, the FSM must leave idle
            within a bounded number of cycles unless reset asserts. &
            \Cref{sec:rtl-controller-mainfsm} \\ \midrule
        \textbf{\emph{op\_done} X-check.} In the four sub-states that
            depend on \emph{op\_done}, the handshake line must not be
            observed at X. Strengthen to a stability/toggle property
            once analog bounds become available. &
            \Cref{sec:rtl-controller-prog,sec:rtl-controller-erase}
        \\ \midrule
        \textbf{Pulse-counter discipline.} When the pulse-mode bus
            carries a non-idle encoding, the pulse counter must stay
            within the valid range until \emph{pulse\_done} asserts
            (one assertion per active command state). &
            \Cref{sec:rtl-controller-pulses} \\ \midrule
        \textbf{One-hot tail discipline.} At the host boundary, the
            structural one-hot validity is a precondition for
            \emph{valid}; mirrored as a constraint in random-stimulus
            environments or a cover in directed ones. &
            \Cref{sec:rtl-pi} \\ \midrule
        \textbf{Recovery branch coverage.} Every regression run must
            observe each of the three verify-check outcomes (match,
            under, over) at least once; runs that exercise the retry
            budget should also observe the retry-exhausted branch. &
            \Cref{sec:rtl-controller-verify-check} \\
        \bottomrule
    \end{tabularx}
\end{table}

The integration testbench and the program--verify regression both
already instantiate the controller and reference its internal state for
the \emph{op\_done} generator, so attaching an assertion binding
introduces no new plumbing.

\section{Future Work --- Constrained-Random Stimulus and Coverage}
\label{sec:future-cr}

The existing environment already exposes a legal-packet generator: the
helper functions that build a transaction word from indices, and the
one-hot validity predicates, together form a constructive way to
generate and filter random host packets. A natural extension is a
constrained-random environment with four pieces. First, randomize host
data and command subject to one-hot tail validity and the enumerated
command set. Second, constrain to the protocol: the controller does not
support pipelined host issue, so the random stream must idle until busy
drops, mirroring the directed bench discipline. Third, define
covergroups on the main state, the three sub-FSM states, and the three
verify-check branches, together with crosses of command class at idle.
Fourth, reuse the shadow-matrix scoreboard from the existing
program--verify regression to self-check each randomly generated
transaction.

\section{Future Work --- Recovery Corner Cases and Gate-Level Flow}
\label{sec:future-next}

\paragraph{Parameterized retry exhaustion.}
A focused bench that holds the under-programmed injection for the
entire duration of a single transaction (rather than just the first
verify cycle) would exhaust the per-cell retry budget on a single cell
and exercise the erase-retry-exhausted branch without programming
hundreds of weights.

\paragraph{Host vs.\ verify erase entry.}
A bench that drives both entry paths into the erase sub-FSM in the
same run, observing the verify/host erase flag register through
hierarchical access, would directly cover the host/verify distinction
that is currently only inferred from end-of-run scoreboards.

\paragraph{Gate-level flow.}
Once the RTL and functional verification stabilize, synthesis, static
timing analysis, and power/area characterization are the natural next
step. Those require a target library, clock/reset constraints, and a
top-level integrating module, none of which are in scope here. The
recovery-policy and integration regressions both remain valid
post-synthesis check-ins, so the functional verification established
here can be reused as a gate-level confidence baseline.
